\section{Rumusan Masalah}

Berdasarkan kesenjangan penelitian yang teridentifikasi, rumusan masalah yang ada pada penelitian ini disusun sebagai berikut:
\begin{enumerate}
    \item Mekanisme pemodelan probabilistik dan pembelajaran struktur dependensi seperti apa yang memungkinkan korelasi antar-atribut ABAC ditangkap tanpa menyimpan populasi solusi eksplisit, sehingga jejak memorinya tidak tumbuh terhadap ukuran populasi?
    \item Bagaimana kualitas kebijakan ABAC yang dihasilkan algoritma usulan dibandingkan terhadap algoritma berpopulasi (BOA, ACO, UBK) dan \textit{compact} murni (\textit{compact-GA/PSO}) pada berbagai tingkat korelasi atribut, dan \textit{trade-off} apa yang muncul antara kualitas kebijakan dan kehematan sumber daya?
    \item Sejauh mana pemisahan frekuensi pembaruan struktur dependensi dari pembaruan parameter, termasuk strategi \textit{warm-start} setelah \textit{drift}, memengaruhi waktu pembaruan dan kualitas kebijakan pada berbagai tingkat \textit{drift} data akses, dibandingkan pembaruan struktur setiap siklus dan penambangan ulang dari awal?
    \item Bagaimana jejak memori algoritma usulan berperilaku terhadap ukuran populasi virtual (NP) dan dimensi ruang solusi (D) ketika dijalankan pada perangkat \textit{edge} target, dan bagaimana kelayakan sumber dayanya dibandingkan algoritma berpopulasi eksplisit (BOA, ACO, UBK)?
\end{enumerate}
Ketersediaan dataset yang memenuhi karakteristik pada rumusan masalah di atas dibahas dan dikonstruksi pada metodologi penelitian.